\section{Results and Analysis Plan}
\label{sec:results}
\draftnote{No pretrained checkpoint from the three-model protocol has been
evaluated for this study.
The tables below are reporting layouts and contain no measured model
performance.}

\begin{table*}[t]
\centering
\small
\setlength{\tabcolsep}{4pt}
\begin{tabular}{lccccccc}
\toprule
 & \multicolumn{4}{c}{Condition accuracy $\uparrow$}
 & \multicolumn{3}{c}{Grounding diagnostic}\\
\cmidrule(lr){2-5}\cmidrule(lr){6-8}
Method & $O^+C^+$ & $O^-C^+$ & $O^+C^-$ & $O^-C^-$
 & All four $\uparrow$ & FPR $\downarrow$
 & Achieved TPR $\uparrow$\\
\midrule
Base & \pending & \pending & \pending & \pending & \pending & \pending & \pending\\
Base + threshold & \pending & \pending & \pending & \pending & \pending & \pending & \pending\\
VTI (adapted) & \pending & \pending & \pending & \pending & \pending & \pending & \pending\\
GrAInS (adapted) & \pending & \pending & \pending & \pending & \pending & \pending & \pending\\
REVIS (adapted) & \pending & \pending & \pending & \pending & \pending & \pending & \pending\\
RUDDER (adapted) & \pending & \pending & \pending & \pending & \pending & \pending & \pending\\
HONES (pending) & \pending & \pending & \pending & \pending & \pending & \pending & \pending\\
Random heads + gain fit & \pending & \pending & \pending & \pending & \pending & \pending & \pending\\
\method & \pending & \pending & \pending & \pending & \pending & \pending & \pending\\
\bottomrule
\end{tabular}
\caption{Planned forced-choice comparison on Qwen2.5-VL-7B. $O$ denotes target
presence and $C$ associated-context presence. All four requires every
condition in a complete original-image group to be correct. Condition and
group accuracy use the zero margin threshold, except for the calibrated
Base + threshold control. FPR and achieved TPR use separately
development-selected operating points for every method. Generated-answer
and caption metrics are reported separately. No entries have been measured.
Additional close comparators in Sec.~\ref{sec:related} remain necessary
before finalizing the table.}
\label{tab:main}
\end{table*}

\begin{table*}[t]
\centering
\small
\setlength{\tabcolsep}{4pt}
\begin{tabular}{llccccc}
\toprule
Model & Method & POPE F1 $\uparrow$ & CHAIR$_i$ $\downarrow$
 & Obj. recall $\uparrow$ & MME $\uparrow$ & MMHal $\uparrow$\\
\midrule
Qwen2.5-VL-7B & Base & \pending & \pending & \pending & \pending & \pending\\
 & Best adapted baseline & \pending & \pending & \pending & \pending & \pending\\
 & \method & \pending & \pending & \pending & \pending & \pending\\
Qwen3-VL-8B & Base & \pending & \pending & \pending & \pending & \pending\\
 & Best adapted baseline & \pending & \pending & \pending & \pending & \pending\\
 & \method & \pending & \pending & \pending & \pending & \pending\\
Gemma 4 E4B & Base & \pending & \pending & \pending & \pending & \pending\\
 & Best adapted baseline & \pending & \pending & \pending & \pending & \pending\\
 & \method & \pending & \pending & \pending & \pending & \pending\\
\bottomrule
\end{tabular}
\caption{Planned three-model transfer comparison. The strongest adapted
baseline is selected separately on development data for each model; the
full method-wise results and output lengths will also be reported. Model
capacities differ, so conclusions concern within-model changes from Base.
All entries are unmeasured.}
\label{tab:transfer}
\end{table*}

\subsection{Implementation verification}
The implementation has passed 74 local CPU correctness tests and static
code checks. Tests include tensor interventions, gradient flow to bounded
gains, restoration of hooks and base weights, data-group separation, and
evaluation adapters. Small randomly initialized instances of the actual
Qwen2.5-VL, Qwen3-VL, Gemma 4, and LLaVA architecture classes exercise image inputs,
grouped-query head slicing, and cached generation. These tests establish
software behavior under their test conditions. They do not reproduce
published baseline scores, establish usefulness of the selected heads, or
predict pretrained-model performance.

\subsection{Does the intervention improve grounding?}
Table~\ref{tab:main} is designed to expose whether fewer false positives
come at the expense of real-object recognition. The intended supporting
pattern is improved performance in the target-absent, associated-context
condition without deterioration when a target is present without its usual
context. A stronger claim requires improvement in complete-group
correctness and in discrimination beyond the base threshold control.
If only the operating point changes, we will report that narrower outcome.

\subsection{Does causal selection matter?}
The crucial comparison holds gain fitting fixed and changes only head
selection. We will inspect whether control-adjusted selection outperforms
random and magnitude-based choices across discovery samples. We will also
compare localization scores with held-out intervention effects, rather
than infer the latter from patch recovery. Unstable selected sets, uniformly
nonpositive scores, or no advantage over random heads would weaken the
proposed mechanism.

\begin{table}[t]
\centering
\small
\setlength{\tabcolsep}{2.5pt}
\begin{tabular}{lcccc}
\toprule
Variant & CHAIR$_i$ $\downarrow$ & Rec.\ $\uparrow$
 & Length & All four $\uparrow$\\
\midrule
Base & \pending & \pending & \pending & \pending\\
Best baseline & \pending & \pending & \pending & \pending\\
\method & \pending & \pending & \pending & \pending\\
No KL & \pending & \pending & \pending & \pending\\
No control penalty & \pending & \pending & \pending & \pending\\
Random heads & \pending & \pending & \pending & \pending\\
\bottomrule
\end{tabular}
\caption{Planned preservation and selection ablations. Rec.\ denotes object
recall. The strongest steering baseline is selected using development
data. Actual caption length is reported because an identical maximum token
budget does not ensure equally informative outputs. All values are pending.}
\label{tab:ablation}
\end{table}

\subsection{Does information preservation transfer?}
Table~\ref{tab:ablation} tests whether first-token regularization transfers
to descriptions. Lower CHAIR with substantially shorter captions or lower
object recall would not demonstrate the intended benefit. We will inspect
hallucination--recall curves and MME changes across all three models, and
report mixed or negative outcomes explicitly. Semantic-context, natural-image,
and capability results should support compatible conclusions before
describing \method as context-robust.
